% 布尔合取的删除式学习（第5章5.4.3）。p=3，目标合取 x1 ∧ ¬x3（学习器不知道）。
% 样本：(1,0,0)->1，(0,1,0)->0，(1,1,0)->1，(1,1,1)->0（均由目标合取标注）。
% 每读到正例，删去在该输入上取 0 的文字；负例不改动文字，只用于核对。
% 编码：蓝底=保留；红底带斜线=本步删除；灰色空格=此前已删除。
% 步1删去 ¬x1、x2、x3；步3删去 ¬x2；最终保留 x1、¬x3。
\begingroup
\begin{tikzpicture}[figure picture,x=1mm,y=1mm,
  keep/.style={draw=FigureBlue,fill=FigureBlueFill,line width=.6pt,minimum width=9mm,minimum height=6mm,inner sep=0pt},
  drop/.style={draw=FigureRed,fill=FigureRedFill,line width=.6pt,minimum width=9mm,minimum height=6mm,inner sep=0pt},
  gone/.style={draw=FigureGrid,fill=white,line width=.6pt,minimum width=9mm,minimum height=6mm,inner sep=0pt,text=FigureGrid}]
\foreach \lit [count=\j] in {x_1,\neg x_1,x_2,\neg x_2,x_3,\neg x_3}{
  \node[figure heading] at (30+10*\j,42) {$\lit$};}
\node[figure heading] at (3,42) {步骤};
\node[figure heading] at (19,42) {输入，标签};
\node[figure heading] at (113,42) {当前合取};
% state codes: k keep, d dropped now, g gone earlier
\foreach \row/\step/\inp/\codes/\conj in {
  0/初始/{全部 $6$ 个文字}/{k,k,k,k,k,k}/{恒为 $0$},
  1/1/{$(1,0,0)$，$1$}/{k,d,d,k,d,k}/{$x_1\wedge\neg x_2\wedge\neg x_3$},
  2/2/{$(0,1,0)$，$0$}/{k,g,g,k,g,k}/{不变；预测 $0$ 正确},
  3/3/{$(1,1,0)$，$1$}/{k,g,g,d,g,k}/{$x_1\wedge\neg x_3$},
  4/4/{$(1,1,1)$，$0$}/{k,g,g,g,g,k}/{不变；预测 $0$ 正确}}{
  \pgfmathsetmacro{\y}{34-8*\row}
  \node at (3,\y) {\step};
  \node at (19,\y) {\inp};
  \foreach \c [count=\j] in \codes {
    \if k\c \node[keep] at (30+10*\j,\y) {};\fi
    \if d\c \node[drop] (dd) at (30+10*\j,\y) {};\draw[FigureRed,line width=.75pt] (dd.south west)--(dd.north east);\fi
    \if g\c \node[gone] at (30+10*\j,\y) {};\fi}
  \node at (113,\y) {\conj};}
\node[figure note,anchor=west] at (0,-6) {蓝：保留；红斜线：本步删除；灰空格：此前已删除。负例只用于核对，不改动文字。};
\end{tikzpicture}
\endgroup
